\chapter{KẾT QUẢ THỰC HIỆN VÀ KẾT LUẬN}\label{ch:ketqua}
\section{Kết quả nghiên cứu và triển khai hệ thống}
\subsection{Tích hợp phần cứng thiết bị đeo và mạng cảm biến cơ thể}
Hệ thống thiết bị đeo được xây dựng dựa trên vi điều khiển trung tâm ESP32. Khối cảm biến tại vùng thắt lưng tích hợp cảm biến quán tính GY-85 và khí áp kế BMP580. Trong đó, mô-đun GY-85 thu thập dữ liệu gia tốc, vận tốc góc và từ trường. Các thông số này đóng vai trò đầu vào cho thuật toán nhận dạng hành động và xác định phương hướng. Cảm biến BMP580 ghi nhận áp suất khí quyển để ước tính biến thiên độ cao giữa các tầng. Toàn bộ tín hiệu thô được xử lý sơ bộ tại biên trước khi truyền sang nút điều khiển trên mũ bảo hộ qua kết nối BLE.

Từ nút thắt lưng, dữ liệu chuyển động được hợp nhất với tín hiệu từ cảm biến khí MQ-9 và MQ-135 gắn trên mũ. Nút trung tâm trên mũ đóng gói toàn bộ thông tin và phát về trạm thu thập qua sóng vô tuyến RF 433 MHz. Cấu trúc phân tầng này giúp tách biệt khối cảm biến đeo trên người với khối truyền thông diện rộng. Đồng thời, hệ thống vẫn duy trì mã định danh đồng nhất cho từng chiến sĩ trong suốt quá trình truyền dẫn dữ liệu.

\subsection{Triển khai mô hình học máy nhận dạng hành động tại biên}
Sau khi thu thập dữ liệu phần cứng, vi điều khiển ESP32 thực hiện quy trình nhận dạng hành động theo thời gian thực. Chuỗi xử lý bao gồm lấy mẫu, lọc nhiễu, phân đoạn cửa sổ trượt, trích xuất đặc trưng và suy luận bằng mô hình Rừng ngẫu nhiên (Random Forest). Tín hiệu gia tốc ba trục được lấy mẫu ở tần số 50 Hz. Mỗi cửa sổ trượt có độ dài 3 giây, tương ứng với 150 mẫu dữ liệu trên từng trục. Mô hình sử dụng 13 loại đặc trưng miền thời gian với cấu hình 17 cây quyết định và độ sâu tối đa là 14. Cấu hình tối ưu này cho phép thiết bị phân loại chính xác 11 trạng thái vận động đặc thù của lính cứu hỏa.

Kết quả phân loại hành động được đồng bộ với bộ đếm bước, góc định hướng và nồng độ khí độc để tạo thành gói tin giám sát hoàn chỉnh. Hệ thống tách biệt rõ chu kỳ lấy mẫu cảm biến và chu kỳ xuất bản tin cảnh báo. Để đảm bảo độ chính xác khi chạy nhúng, quy trình suy luận trên vi điều khiển tuân thủ nghiêm ngặt đơn vị đo, thứ tự đặc trưng và bảng gán nhãn từ giai đoạn huấn luyện.

\subsection{Ước lượng phương hướng và định vị tương đối trong nhà}
Song song với nhánh nhận dạng hành động, dữ liệu quán tính tiếp tục được khai thác cho bài toán định vị tương đối trong nhà (PDR). Thuật toán sử dụng độ lớn gia tốc tổng hợp để phát hiện pha bước chân. Tọa độ của người lính được cập nhật liên tục dựa trên độ dài sải bước và góc hướng tại thời điểm bước chân hoàn thành. Hệ trục tọa độ mặt phẳng được quy ước chuẩn hóa với trục $x$ hướng Đông và trục $y$ hướng Bắc. Góc định hướng lấy mốc $0^\circ$ tại hướng Bắc và tăng dần theo chiều kim đồng hồ. Vị trí mới được tính toán tích lũy theo từng bước đi, trong khi tổng quãng đường di chuyển được ghi nhận để phục vụ thống kê thực nghiệm.

Trong thuật toán cập nhật tọa độ, hệ thống sử dụng giá trị góc liên tục để đảm bảo độ chính xác hình học. Khi người lính xoay người tại chỗ mà không bước đi, hệ thống chỉ cập nhật góc hướng và giữ nguyên tọa độ vị trí. Kết quả phân tích tại Chương~\ref{ch:kiemthu} chỉ ra rằng hiện tượng trôi góc tích lũy tác động trực tiếp đến sai số quỹ đạo. Do đó, việc hiệu chuẩn cảm biến và kiểm soát chất lượng góc quay là yêu cầu cốt lõi của phân hệ định vị.

\subsection{Thiết kế trạm thu thập dữ liệu dựa trên Raspberry Pi 3B và mô-đun SIM 4G}
Để đưa dữ liệu từ hiện trường về sở chỉ huy, hệ thống sử dụng trạm thu thập trung gian (Gateway) đặt bên ngoài khu vực nguy hiểm. Trạm Gateway được thiết kế dựa trên máy tính nhúng Raspberry Pi 3 Model B kết hợp với bộ thu RF 433 MHz và mô-đun SIM 4G. Hệ thống thực hiện tuần tự các tác vụ: tiếp nhận khung truyền RF từ thiết bị đeo, kiểm tra mã lỗi, chuẩn hóa cấu trúc dữ liệu và đẩy bản tin lên máy chủ MQTT Broker qua mạng di động 4G. Các trường thông tin gồm mã định danh thiết bị, số thứ tự gói tin và nhãn thời gian được giữ nguyên vẹn để phục vụ đối soát dữ liệu đầu cuối.

Phần mềm trên Raspberry Pi tích hợp cơ chế hàng đợi bộ đệm nhằm lưu trữ tạm thời các bản tin khi sóng 4G bị gián đoạn. Khi kết nối mạng phục hồi, thuật toán tự động kết nối lại và đồng bộ dữ liệu kết hợp chống trùng lặp gói tin. Bản tin giám sát cũng bao gồm cờ trạng thái kết nối và dấu thời gian của mẫu đo cuối cùng. Thiết kế này giúp hệ thống phân biệt rõ dữ liệu thời gian thực và dữ liệu được truyền bù. Các chỉ tiêu kỹ thuật về tỷ lệ mất gói, độ trễ truyền dẫn và khả năng tự khôi phục được kiểm chứng chi tiết tại Chương~\ref{ch:kiemthu}.


\subsection{Minh chứng thiết kế phần cứng và nhận bản tin MQTT}
Sơ đồ nguyên lý ESP32-C5--GY-85 cùng mô hình bố trí linh kiện PCB hai mặt được trình bày tại Hình~\ref{fig:pcb-schematic}, Hình~\ref{fig:pcb-3d-tren} và Hình~\ref{fig:pcb-3d-duoi}. Các tư liệu này ghi nhận kết quả thiết kế bo mạch ở mức bản vẽ và mô phỏng 3D, chưa thay thế quy trình lắp ráp và đo kiểm phần cứng thực tế. Ảnh chụp tại Hình~\ref{fig:mqttx-minhchung} minh chứng việc MQTTX nhận được bản tin JSON chứa trạng thái hành động, thông số bước chân, thông tin hướng và chuỗi mẫu IMU/PDR. Việc đánh giá chất lượng đường truyền RF--4G đầu cuối cần những phép đo độc lập bổ sung.

\section{Tổng hợp các thành phần đã hoàn thiện}
\begingroup\SmallTable
\begin{longtable}{|P{106bp}|P{147bp}|P{155bp}|}
\caption{Tổng hợp nội dung phần cứng, xử lý và truyền dữ liệu}\label{tab:completion}\\\hline
\textbf{Hạng mục}&\textbf{Nội dung}&\textbf{Vai trò trong hệ thống}\\\hline\endfirsthead
\hline\textbf{Hạng mục}&\textbf{Nội dung}&\textbf{Vai trò trong hệ thống}\\\hline\endhead
Thiết bị đeo&Vi điều khiển ESP32, cảm biến GY-85 và BMP580; lấy mẫu, xử lý và đóng gói dữ liệu.&Thu thập thông số vận động và biến thiên áp suất khí quyển tại hiện trường.\\\hline
Bo mạch thử nghiệm&Sơ đồ nguyên lý ESP32-C5--GY-85; mô hình 3D hai mặt PCB.&Tài liệu phục vụ đối chiếu đấu nối và bố trí linh kiện trước chế tạo.\\\hline
Nhận dạng hành động&Cửa sổ 150 mẫu (3 giây); 13 đặc trưng miền thời gian; mô hình Random Forest 17 cây.&Phân loại 11 trạng thái vận động theo thời gian thực ngay trên vi điều khiển.\\\hline
Định vị tương đối (PDR)&Phát hiện bước chân, ước lượng góc hướng và cập nhật tọa độ theo từng chu kỳ bước.&Tái lập quỹ đạo di chuyển tương đối của chiến sĩ từ vị trí xuất phát.\\\hline
Trạm Gateway&Máy tính nhúng Raspberry Pi 3B, bộ thu RF 433 MHz, mô-đun SIM 4G và giao thức MQTT.&Thu thập gói tin RF, lưu đệm chống mất mẫu và chuyển tiếp dữ liệu lên máy chủ đám mây.\\\hline
Tiếp nhận MQTT&Ảnh chụp MQTTX trên chủ đề \nolinkurl{ssalab/firefighter1}, bản tin dữ liệu dạng JSON.&Minh chứng nhận dữ liệu ở phía MQTT client trong một phiên kiểm tra.\\\hline
\end{longtable}\endgroup

\section{Hạn chế tồn tại và nguyên nhân kỹ thuật}\label{sec:limitations}
\begin{itemize}
\item \textbf{Hiện tượng trôi góc và sai lệch phương hướng:} Sai số lệch chuẩn (bias) của con quay hồi chuyển tích lũy tăng dần theo thời gian vận hành. Bên cạnh đó, nhiễu từ trường phát sinh từ pin, đường dây cấp nguồn và các kết cấu kim loại trong tòa nhà làm suy giảm độ chính xác của từ kế. Sự xê dịch tư thế gắn thiết bị trên cơ thể người lính cũng gây ra sai số khi ước lượng hướng di chuyển.
\item \textbf{Sai số tích lũy trong định vị PDR:} Việc phát hiện thiếu hoặc thừa bước chân tác động trực tiếp đến kết quả tính quãng đường. Độ dài sải bước thực tế luôn biến thiên theo tốc độ và thể trạng chứ không cố định. Ngoài ra, sự lệch pha giữa hướng của thân người và hướng di chuyển thực tế làm sai số vị trí tăng dần theo chiều dài quỹ đạo. Mô hình đếm bước hiện tại mới tối ưu cho trạng thái đi bộ và còn hạn chế khi người lính bò, trườn hoặc leo trèo.
\item \textbf{Giới hạn về phần cứng và độ trễ xử lý:} Hiện tượng sụt áp nguồn, lỏng giắc kết nối hoặc rung lắc cơ học có thể gây mất mẫu dữ liệu và làm khởi động lại vi điều khiển. Các cảm biến khí bán dẫn MQ đòi hỏi thời gian làm nóng màng nhạy và phải hiệu chuẩn định kỳ. Đồng thời, việc sử dụng cửa sổ thu thập dữ liệu cố định 3 giây cũng tạo ra một độ trễ nhất định trong việc phát hiện sự cố tức thời.
\item \textbf{Nguy cơ gián đoạn đường truyền vô tuyến:} Sóng RF 433 MHz dễ bị suy hao khi xuyên qua nhiều lớp tường bê tông dày hoặc xảy ra xung đột kênh truyền. Tương tự, kết nối từ mô-đun SIM 4G phụ thuộc lớn vào chất lượng phủ sóng di động và lưu lượng mạng tại khu vực cháy. Nếu thiếu cơ chế đồng bộ thời gian và lọc trùng lặp hiệu quả, các gói tin truyền bù có thể gây sai lệch khi diễn giải trạng thái hiện tại của thiết bị.
\end{itemize}

\section{Định hướng khắc phục và hoàn thiện hệ thống}\label{sec:improvement}
\begin{itemize}
\item \textbf{Hoàn thiện thuật toán hiệu chuẩn và ước lượng hướng:} Hệ thống cần tự động bù sai số tĩnh của con quay hồi chuyển và hiệu chuẩn từ kế ngay trên bo mạch hoàn chỉnh. Thuật toán hợp nhất cảm biến phải tích hợp bước kiểm tra nhiễu từ trường và bù góc nghiêng cơ thể trước khi tính toán góc hướng. Khi phát hiện từ trường biến động bất thường, thuật toán cần tự động hạ trọng số của từ kế để bảo toàn độ tin cậy.
\item \textbf{Nâng cao độ chính xác của mô hình định vị PDR:} Độ dài bước chân cần được ước lượng thích nghi theo từng cá nhân và vận tốc di chuyển thực tế. Thuật toán phát hiện bước cần tự động điều chỉnh ngưỡng theo nhãn hành động nhận dạng được. Ngoài ra, hệ thống có thể kết hợp thêm các mốc tham chiếu vị trí hoặc áp dụng thuật toán khớp bản đồ (map matching) để triệt tiêu sai số tích lũy.
\item \textbf{Tối ưu hóa thiết kế phần cứng và chương trình nhúng:} Cảm biến quán tính GY-85 cần được cố định chắc chắn vào khung vỏ để giảm nhiễu cơ học. Mạch nguồn phải được phân tách độc lập giữa khối vi điều khiển và khối công suất RF, mô-đun SIM 4G. Trên vi điều khiển, tác vụ lấy mẫu cảm biến cần được tách biệt khỏi tác vụ truyền thông qua cơ chế đa nhiệm. Độ dài cửa sổ trượt và bước trượt cũng cần được điều chỉnh linh hoạt dựa trên chỉ số phát hiện bất thường.
\item \textbf{Tăng cường độ tin cậy của giao thức truyền thông:} Trạm Gateway sử dụng Raspberry Pi 3B và mô-đun SIM 4G cần quản lý chặt chẽ bộ đệm hàng đợi, số thứ tự gói tin và nhãn thời gian gốc. Cơ chế tự động kết nối lại và truyền bù dữ liệu phải hoạt động ổn định khi mạng chập chờn. Cần mở rộng kiểm thử hệ thống trong điều kiện mất kết nối diện rộng, khởi động lại đột ngột và hoạt động đồng thời nhiều nút mạng. 
\end{itemize}

\section{Kết luận chung}
Nhóm đã thiết kế và triển khai thành công kiến trúc hệ thống giám sát an toàn cho lính cứu hỏa theo thời gian thực. Hệ thống hoàn thiện bao gồm thiết bị đeo tại thắt lưng (ESP32, GY-85, BMP580), nút tổng hợp gắn trên mũ bảo hộ, mạng truyền thông hỗn hợp BLE/RF và trạm Gateway sử dụng Raspberry Pi 3B kết hợp mô-đun SIM 4G. Các thuật toán nhận dạng hành động học máy tại biên và định vị quán tính PDR đã được tích hợp đồng bộ vào quy trình xử lý, đóng gói và chuyển tiếp dữ liệu.

Các kết quả thực nghiệm bước đầu khẳng định tính khả thi của giải pháp trong việc theo dõi trạng thái vận động, cảnh báo môi trường và ước lượng vị trí tương đối bên trong tòa nhà. Hướng phát triển tiếp theo sẽ tập trung vào việc giảm thiểu sai số trôi góc, tối ưu hóa mô hình bước chân thích nghi, nâng cao độ bền phần cứng và kiểm chứng độ tin cậy truyền thông đầu cuối trong các điều kiện diễn tập sát với thực tế.